% ============================================================
\chapter{Modul 12: Jupyter Notebooks}
% ============================================================

\section{Lektion 34: Arbeiten wie ein Ingenieur}

Jupyter Notebooks kombinieren Code, Ergebnisse, Diagramme und Erklärungen in einem einzigen Dokument. Für Robotik-Ingenieure sind sie das ideale Werkzeug zur Datenanalyse und zum Experimentieren.

\subsection*{Installation und Start}

\begin{lstlisting}
# Terminal/Kommandozeile:
pip install jupyterlab
jupyter lab   # Startet im Browser
\end{lstlisting}

\subsection*{Notebook-Zellen}

Ein Notebook besteht aus \textbf{Zellen}. Jede Zelle ist entweder Code oder Markdown-Text.

\begin{lstlisting}
# Code-Zelle: wird ausgefuehrt mit Shift+Enter
import numpy as np
import matplotlib.pyplot as plt

# Ergebnis erscheint direkt unter der Zelle
x = np.linspace(0, 2*np.pi, 100)
y = np.sin(x)
print(f"Max: {y.max():.4f}")
\end{lstlisting}

\begin{merksatz}
Notebooks sind \textbf{nicht} für Produktions-Code -- sie sind für Analyse und Experimente. Den fertigen, getesteten Code lagert man in \texttt{.py}-Dateien aus, die das Notebook importiert.
\end{merksatz}

\section{Lektion 35: Sensordaten analysieren}

\begin{lstlisting}
import numpy as np
import pandas as pd   # pip install pandas

# CSV-Sensorlog laden
df = pd.read_csv("sensor_log.csv")

# Schnelle Uebersicht
print(df.describe())     # Statistik aller Spalten
print(df.head())         # Erste 5 Zeilen

# Filtern: nur Messungen mit Hindernis
hindernisse = df[df["abstand_m"] < 0.5]
print(f"Hindernisereignisse: {len(hindernisse)}")

# Zeitreihe glaetten (gleitender Mittelwert)
df["abstand_geglaettet"] = df["abstand_m"].rolling(window=5).mean()

# Exportieren
df.to_csv("analysiert.csv", index=False)
\end{lstlisting}

\begin{praxis}
Nach einem Testlauf mit dem echten Roboter öffnet man das Notebook, lädt den Log und analysiert in wenigen Minuten: Wo war die Fehlerrate am höchsten? Wie hat sich der Akku verhalten? Welche Sensorwerte waren Ausreißer?
\end{praxis}

\begin{uebung}[title=Projektaufgabe Modul 12: Testfahrt-Analyse]
Generiere einen simulierten Fahrt-Log (100 Zeilen, CSV: Zeit, X, Y, Akku, Abstand\_vorne). Lade ihn in ein Jupyter Notebook. Berechne: Gesamtstrecke (aus X/Y-Positionen), Akku-Verbrauch pro Meter, Anteil der Zeit mit Hindernis $<$ 0.5 m.
\end{uebung}
